\chapter{System Requirements Specification}

\section{Introduction}

This chapter presents a comprehensive and detailed analysis of the system requirements for the E-Commerce Microservices Platform with Event-Driven Architecture. The requirements specification serves as the foundation for the entire system design and implementation, clearly defining what the system must accomplish, how it should behave under various conditions, and what constraints must be respected during development. The requirements have been systematically gathered through extensive analysis of industry best practices in e-commerce platform development, careful examination of existing microservices implementations, thorough study of distributed systems literature, and identification of technical necessities specific to cloud-native architecture patterns.

The requirements are categorized into functional requirements, which specify the behaviors and capabilities the system must provide to satisfy user needs and business objectives, and non-functional requirements, which define the quality attributes, performance characteristics, and operational constraints that the system must satisfy. Each requirement is uniquely identified, clearly stated, and traceable to specific system components and user needs.

\section{Functional Requirements}

Functional requirements describe the specific behaviors, operations, and capabilities that the system must exhibit. These requirements define what the system should do when presented with particular inputs or when operating under specific conditions. The functional requirements for this e-commerce platform are organized by service domain, reflecting the microservices decomposition strategy.

\subsection{User Management and Authentication Requirements}

The User Service serves as the foundation for identity management and access control across the entire platform. It must provide comprehensive user lifecycle management capabilities that support secure authentication, authorization, and profile management while integrating seamlessly with the API Gateway for token-based access control.

The system shall enable new users to create accounts through a registration process that collects essential information including email address, password, first name, last name, and optionally a phone number. During registration, the system must perform comprehensive validation of all input fields, ensuring that email addresses conform to standard format specifications and do not already exist in the system to prevent duplicate accounts. Password validation must enforce complexity requirements mandating minimum length of eight characters, inclusion of uppercase and lowercase letters, numerical digits, and special characters to ensure strong credential security. Upon successful registration, the system shall store the user's password using bcrypt hashing algorithm with a cost factor of twelve, providing robust protection against brute force and rainbow table attacks.

Authentication shall be implemented through a secure login process where users provide their email address and password credentials. The system shall retrieve the user record from the database based on the provided email address, extract the stored bcrypt hash, and verify the provided password against this hash using bcrypt's verification mechanism. Upon successful authentication, the system shall generate a JSON Web Token (JWT) containing the user's unique identifier, email address, first name, last name, and assigned roles encoded in the token payload. The token shall be signed using HMAC-SHA256 algorithm with a secure secret key and configured with an expiration time of twenty-four hours from the issuance timestamp. This token shall be returned to the client and must be included in the Authorization header of all subsequent API requests to protected endpoints.

The system shall implement role-based access control with support for multiple user roles including standard USER role and elevated ADMIN role. During user registration, new accounts shall automatically be assigned the USER role by default, granting access to standard e-commerce functionality such as browsing products, managing shopping carts, and placing orders. The ADMIN role shall be explicitly assigned by existing administrators and grants additional privileges including product catalog management, order administration, user account management, and access to administrative dashboards. The API Gateway shall validate JWT tokens on incoming requests, extract the user's roles from the token payload, and enforce authorization policies by rejecting requests to protected endpoints when the user lacks the required role.

Profile management functionality shall allow authenticated users to update their personal information including first name, last name, phone number, and password. When updating passwords, the system shall require the user to provide their current password for verification purposes before applying the new password, preventing unauthorized password changes in scenarios where sessions are left open on shared devices. All profile updates shall be validated for data integrity and business rule compliance before being persisted to the database.

\subsection{Product Catalog Management Requirements}

The Product Service is responsible for comprehensive management of the e-commerce product catalog, providing capabilities for product creation, retrieval, search, filtering, updates, and deletion while maintaining data integrity and supporting high-performance queries for customer-facing operations.

Product creation shall be restricted to users with ADMIN role privileges to maintain catalog quality and prevent unauthorized product listings. When creating a new product, administrators must provide the product name, description, price, and category as mandatory fields, while stock quantity, image URL, and active status are optional with sensible defaults. The system shall validate that the product name is non-empty and does not exceed two hundred characters, the price is a positive decimal value with maximum two decimal places representing currency units, the category is selected from a predefined list of valid categories, and the stock quantity is a non-negative integer. Upon successful validation, the system shall persist the product to the PostgreSQL database, automatically generating a UUID as the unique product identifier and recording the creation timestamp using Spring Data JPA auditing capabilities.

Product retrieval shall support multiple query patterns to accommodate different use cases. Individual product retrieval by unique identifier shall return complete product details including all attributes and current stock levels, enabling product detail pages and order validation workflows. Product listing shall return paginated results with configurable page size defaulting to twenty products per page, including metadata about total element count, total page count, and current page number to support pagination controls in the user interface. The system shall implement database-level pagination using SQL LIMIT and OFFSET clauses to ensure efficient query execution even with large product catalogs containing thousands of items.

Full-text search functionality shall enable customers to find products by searching across product names, descriptions, and categories simultaneously. The system shall leverage PostgreSQL's built-in full-text search capabilities by creating tsvector columns that combine searchable text fields and indexing these columns using Generalized Inverted Index (GIN) for efficient query execution. Search queries shall be converted to tsquery format, matched against indexed tsvectors, and ranked by relevance using the ts\_rank function which considers term frequency and inverse document frequency to surface the most relevant results first. Search results shall be returned in descending order of relevance score with the most relevant products appearing first in the result set.

Product filtering shall allow customers to narrow down product listings based on multiple criteria including category, price range, and availability status. Category filtering shall restrict results to products belonging to the specified category, supporting hierarchical category structures where selecting a parent category includes products from all child categories. Price range filtering shall accept minimum and maximum price parameters and return only products whose prices fall within the specified range, enabling budget-conscious shopping experiences. Availability filtering shall distinguish between active products with positive stock quantities and inactive or out-of-stock products, with the default behavior showing only available products to improve customer experience.

Product updates shall allow administrators to modify any product attribute including name, description, price, category, stock quantity, image URL, and active status. When updating product prices, the system shall apply the new price only to future orders while preserving the historical prices recorded in existing order line items, ensuring accurate financial records and preventing disputes over pricing changes. Stock quantity updates shall coordinate with the Inventory Service through event-driven communication to maintain consistency between catalog information and actual inventory levels. Product deletion shall implement soft deletion by setting the active status to false rather than physically removing database records, preserving referential integrity for historical orders that reference deleted products and enabling recovery of accidentally deleted products.

\subsection{Order Processing and Saga Orchestration Requirements}

The Order Service implements the most complex business logic in the platform, orchestrating the complete order lifecycle through the Saga pattern while coordinating with Inventory and Payment services through event-driven communication to ensure distributed transaction consistency.

Order creation shall accept authenticated requests containing the customer's unique identifier and an array of order items, where each item specifies a product identifier and desired quantity. The system shall validate that the customer exists and is authenticated, each referenced product exists in the catalog, and requested quantities are positive integers. The system shall retrieve current product prices from the Product Service through synchronous REST API calls or cached data, calculate line item totals by multiplying unit prices by quantities, and compute the order grand total by summing all line item totals. Upon successful validation and price calculation, the system shall persist the order to the database with PENDING status, indicating that the order has been created but the Saga workflow has not yet completed.

Following successful order persistence, the system shall publish an OrderCreated event to the Kafka event bus, initiating the Saga workflow that coordinates inventory reservation and payment processing across multiple services. The OrderCreated event shall contain comprehensive order information including order identifier, customer identifier, array of order items with product identifiers and quantities, total amount, and timestamp. This event shall be published using the outbox pattern, where the event is first stored in an outbox table within the same database transaction as the order record, ensuring atomicity between order creation and event generation, and subsequently published to Kafka by a separate outbox processor to guarantee reliable event delivery even in the presence of transient failures.

The Order Service shall consume events published by other services as part of the Saga workflow, advancing the order state machine based on received events. When the Inventory Service successfully reserves inventory for all order items, it shall publish a StockReserved event that the Order Service consumes to advance the order workflow to the payment processing phase. If the Inventory Service cannot reserve inventory for one or more items due to insufficient stock, it shall publish a StockReservationFailed event that triggers the Order Service to cancel the order by transitioning to CANCELLED status and publishing an OrderCancelled event to notify other services.

When the Payment Service successfully processes payment for the order, it shall publish a PaymentCompleted event that the Order Service consumes to transition the order to CONFIRMED status, indicating successful completion of the Saga workflow. The Order Service shall then publish an OrderConfirmed event that triggers the Notification Service to send order confirmation emails to the customer. If the Payment Service fails to process payment due to declined cards, insufficient funds, or gateway errors, it shall publish a PaymentFailed event that triggers compensating transactions including order cancellation and inventory release.

The Order Service shall maintain comprehensive order history tracking all status transitions with timestamps, enabling customers to monitor order progress and administrators to audit order processing. The system shall provide query endpoints for retrieving individual orders by identifier, listing orders filtered by customer, status, and date range, and retrieving order status history showing the complete timeline of status changes. Customer-facing queries shall enforce access control ensuring customers can only view their own orders, while administrative queries shall provide unrestricted access for support and analytics purposes.

\subsection{Payment Processing Requirements}

The Payment Service handles all financial transaction processing with integration to external payment gateways, implementing robust error handling, transaction recording, and refund capabilities required for production e-commerce operations.

Payment processing shall accept requests containing order identifier, payment method specification, payment method details, and transaction amount. The system shall validate that the referenced order exists and is in PENDING or CONFIRMED status awaiting payment, the payment amount matches the order total to prevent underpayment or overpayment, and the payment method details conform to the requirements of the selected payment gateway. For credit card payments via Stripe integration, the system shall accept tokenized card details provided by Stripe Elements or Stripe.js client-side libraries, ensuring that raw card numbers never touch the application servers to maintain PCI DSS compliance. For PayPal payments, the system shall redirect users to PayPal's authorization flow and process the returned authorization codes. For bank transfers, the system shall generate payment instructions and await confirmation of transfer completion.

The system shall maintain a payment state machine tracking each payment through its lifecycle from initiation to completion or failure. When a payment request is first received, the payment shall be created in INITIATED status. As the system communicates with the external payment gateway, the status shall transition to PROCESSING. Upon receiving a successful response from the payment gateway, the status shall transition to COMPLETED and a PaymentCompleted event shall be published to Kafka. If the payment gateway declines the transaction due to insufficient funds, expired cards, fraud detection, or other reasons, the status shall transition to FAILED and a PaymentFailed event shall be published. All state transitions shall be persisted to the database with timestamps to enable audit trails and troubleshooting.

Refund processing shall support both full refunds and partial refunds based on business requirements. Full refunds shall refund the entire payment amount to the original payment method, while partial refunds shall refund a specified amount less than the original payment, useful for scenarios involving returned items or service credits. The refund process shall invoke the payment gateway's refund API with the original transaction identifier and refund amount, update the payment status to REFUNDED upon successful refund, and publish a RefundProcessed event to notify other services of the refund. The system shall record refund transactions separately from original payments to maintain clear financial records and support reconciliation processes.

\subsection{Inventory Management Requirements}

The Inventory Service manages stock levels, processes inventory reservations for orders, and ensures accurate inventory tracking across the platform through event-driven integration with the Order Service.

Stock management shall maintain current inventory quantities for all products in the catalog, initializing stock levels when products are created and adjusting quantities as orders are fulfilled, inventory is received from suppliers, or stock adjustments are made for reasons such as damaged goods or inventory corrections. The system shall enforce non-negative stock constraints preventing inventory quantities from dropping below zero, which would indicate overselling or data integrity issues.

Inventory reservation shall be triggered automatically when the Inventory Service consumes OrderCreated events from Kafka. For each order item, the system shall check if sufficient stock is available by comparing the requested quantity against current available quantity (total stock minus already reserved quantities). If all order items have sufficient available stock, the system shall reserve the quantities by incrementing reserved quantities and decrementing available quantities for each product, then publish a StockReserved event to acknowledge successful reservation. If any order item has insufficient available stock, the system shall not reserve any quantities (implementing all-or-nothing reservation semantics), release any partially reserved quantities, and publish a StockReservationFailed event indicating which products had insufficient stock.

Reservation expiration shall automatically release reserved inventory when orders are not completed within a configured timeout period, typically thirty minutes. This mechanism prevents indefinite inventory holds for abandoned orders, ensuring that stock remains available for other customers. The system shall implement expiration checking through scheduled tasks that query for reservations older than the timeout threshold with associated orders still in PENDING status, releasing these reservations and updating available quantities accordingly.

\section{Non-Functional Requirements}

Non-functional requirements define the quality attributes, performance characteristics, security requirements, and operational constraints that the system must satisfy to meet user expectations and business objectives.

\subsection{Performance Requirements}

The system shall achieve response time targets of less than two hundred milliseconds for simple read operations such as product retrieval by identifier, less than five hundred milliseconds for complex queries involving search and filtering, and less than one second for write operations including order creation and payment processing. These targets shall be measured under normal load conditions representing typical production usage patterns.

The API Gateway shall support a minimum throughput of one thousand requests per second across all services, with individual services capable of handling at least two hundred requests per second for their specific endpoints. The system shall maintain these throughput levels while keeping CPU utilization below seventy percent and memory utilization below eighty percent to provide headroom for traffic spikes.

Kafka message producers shall achieve publish latencies of less than ten milliseconds for ninety-five percent of messages under normal load, ensuring timely event delivery for Saga workflow progression. Kafka consumers shall process messages with end-to-end latency of less than one hundred milliseconds from publication to handler completion, enabling responsive order processing workflows.

\subsection{Scalability Requirements}

The system shall support horizontal scaling of all microservices through Kubernetes Horizontal Pod Autoscaler, automatically increasing replica counts when CPU utilization exceeds seventy percent or custom metrics indicate increased load. The system shall scale from minimum two replicas per service for high availability to maximum ten replicas per service for handling traffic spikes, with scaling decisions evaluated every thirty seconds.

Database connections shall scale efficiently with service replicas, with connection pool sizes automatically adjusted based on the number of service instances. The system shall support up to fifty concurrent database connections per PostgreSQL instance without performance degradation, requiring connection pooling through HikariCP with optimized pool sizing.

Kafka topic partitions shall be configured to support parallel consumer processing, with partition counts sized to accommodate maximum expected consumer parallelism. The system shall support adding partitions to topics without disrupting existing consumers, enabling scalability improvements as consumer counts increase.

\subsection{Reliability and Availability Requirements}

The system shall achieve ninety-nine point nine percent availability for user-facing endpoints, allowing maximum eight point seven six hours of downtime per year for maintenance and unplanned outages. Critical services including API Gateway, User Service, and Order Service shall implement active-active deployments across multiple availability zones to survive zone failures without service disruption.

Circuit breakers shall protect against cascade failures by monitoring failure rates for inter-service calls and opening circuits when failure rates exceed fifty percent over a sliding window of ten requests. Open circuits shall prevent calls to failing services for a wait duration of ten seconds before transitioning to half-open state to test recovery, with half-open state permitting three test calls to determine if the service has recovered.

Data persistence shall ensure durability through database replication and Kafka message replication, with PostgreSQL configured for synchronous replication to at least one standby replica and Kafka configured with replication factor of three for all topics. The system shall survive single node failures without data loss through automatic failover to standby replicas.

\subsection{Security Requirements}

All external communication shall use HTTPS with TLS version one point three or higher, with certificates managed through automated certificate management infrastructure. Internal service-to-service communication shall use HTTP in development environments but shall implement mutual TLS authentication in production environments to prevent unauthorized service access and man-in-the-middle attacks.

Password storage shall use bcrypt hashing algorithm with cost factor of twelve, providing strong protection against brute force attacks while maintaining acceptable authentication performance. JWT tokens shall be signed using HMAC-SHA256 with minimum two hundred fifty-six bit secret keys, with token expiration set to twenty-four hours to limit the window of token misuse if tokens are compromised.

Input validation shall be performed on all user inputs using bean validation annotations and custom validators, preventing SQL injection through parameterized queries, cross-site scripting through output encoding, and injection attacks through strict input validation. The system shall implement rate limiting at the API Gateway level to prevent denial-of-service attacks, with default limits of ten requests per second for authentication endpoints and twenty requests per second for general endpoints.

\subsection{Maintainability Requirements}

Code organization shall follow consistent package structure across all microservices with separate packages for controllers, services, repositories, entities, DTOs, configuration, and utilities, enabling developers to navigate any service codebase efficiently. All services shall use consistent dependency versions managed through parent POM to prevent dependency conflicts and simplify version upgrades.

Logging shall implement structured JSON format with consistent fields including timestamp, log level, service name, trace ID, span ID, thread name, and message, enabling efficient log aggregation and analysis through the ELK Stack. Log levels shall be appropriately configured with INFO level for production operations, DEBUG level for development and troubleshooting, and WARN and ERROR levels for identifying issues requiring attention.

Configuration management shall externalize all environment-specific settings including database connection strings, Kafka broker addresses, Eureka server URLs, and feature flags through Spring Cloud Config or Kubernetes ConfigMaps, enabling configuration changes without code modifications or service restarts for supported properties.

\section{System Constraints}

System constraints represent limitations and restrictions that influence design and implementation decisions, including technology choices, regulatory requirements, and operational limitations.

The system shall be implemented using Java version seventeen or higher to leverage modern language features including records, sealed classes, and pattern matching while maintaining compatibility with Spring Boot three point x framework requirements. All microservices shall use Spring Boot three point one point two as the application framework with Spring Cloud two thousand twenty-two point zero point four for cloud-native patterns, ensuring consistent technology stack across services.

Database technology shall be restricted to PostgreSQL version fifteen or higher for all services, providing ACID compliance, advanced query capabilities, full-text search, and JSON support while simplifying operational management through technology standardization. Message broker technology shall be Apache Kafka version seven point four or higher, providing the event streaming capabilities required for event-driven architecture with guaranteed message ordering and at-least-once delivery semantics.

Container deployment shall use Docker for containerization with multi-stage builds optimizing image sizes, and Kubernetes for orchestration with minimum version one point two four, providing the scalability, self-healing, and deployment automation required for production operations. The system shall operate within resource constraints defined for container deployments, with each service allocated minimum five hundred twelve megabytes of memory and two hundred fifty millicores of CPU, with maximum limits of one gigabyte memory and five hundred millicores CPU to prevent resource exhaustion.

\section{Use Case Specifications}

Use cases provide detailed descriptions of system interactions from user perspectives, capturing the sequence of actions, alternative flows, and exception handling for key system functionalities.

The Place Order use case describes the primary customer journey through the e-commerce platform. The primary actor is a registered customer who has browsed products and added items to their shopping cart. The use case begins when the customer initiates checkout by providing their customer identifier and cart contents. The system validates the customer's authentication status, verifies product availability and current pricing, calculates the order total, creates the order in PENDING status, and publishes the OrderCreated event. The system then displays order confirmation with order identifier and estimated delivery timeline. The alternative flow handles scenarios where products are out of stock, in which case the system informs the customer of unavailable items and suggests similar products. The exception flow handles payment failures, where the system cancels the order, releases reserved inventory, and prompts the customer to try alternative payment methods.

The Manage Product Catalog use case describes administrative functionality for maintaining the product catalog. The primary actor is an administrator user with ADMIN role privileges. The use case includes creating new products with complete attribute specification, updating existing product information including prices and stock levels, searching and filtering products for review, and soft-deleting products that are discontinued. The system validates all inputs against business rules, enforces authorization checks ensuring only administrators can perform these operations, and maintains audit trails of all catalog changes for accountability and troubleshooting.

\section{Requirements Traceability}

Requirements traceability ensures that each requirement is addressed in the design and implementation phases, with clear mapping from requirements to system components and test cases. The functional requirements specified in this chapter directly inform the microservices decomposition strategy, with User Management requirements mapping to the User Service, Product Catalog requirements mapping to the Product Service, Order Processing requirements mapping to the Order Service, Payment Processing requirements mapping to the Payment Service, and Inventory Management requirements mapping to the Inventory Service.

Non-functional requirements influence architectural decisions including the selection of Spring Cloud Gateway for API Gateway pattern implementation to meet performance and scalability requirements, Netflix Eureka for service discovery to support dynamic scaling, Resilience4j for circuit breakers to achieve reliability targets, and comprehensive monitoring infrastructure to enable observability and troubleshooting. Security requirements drive the implementation of JWT authentication, bcrypt password hashing, input validation, rate limiting, and TLS encryption throughout the system.

\section{Summary}

This chapter has presented a comprehensive requirements specification for the E-Commerce Microservices Platform, detailing functional requirements across user management, product catalog, order processing, payment handling, and inventory management domains. Non-functional requirements have been specified for performance, scalability, reliability, security, and maintainability, establishing quality targets that guide architectural decisions. System constraints have been identified regarding technology choices, resource allocations, and operational limitations. Use case specifications have illustrated key system interactions from user perspectives, and requirements traceability has been established linking requirements to system components.

These requirements form the foundation for the system design presented in the following chapter, where architectural decisions, component designs, database schemas, API specifications, and event schemas are developed to satisfy these requirements while adhering to microservices architecture principles and industry best practices.
